\section{Base nonvanishing and removal of the period twist}
\label{sec:nonvanishing}

It remains to prove the section construction stated in the introduction.
The marked quotient gives two cases. A point quotient supplies scalar
coordinates directly; for a positive-dimensional quotient the exact
section lattice of \Cref{sec:relative-iitaka} lets us work on its
intermediate base. There we remove the divisor added by period
positivity. Every argument here is independent of logarithmic
subadditivity.

\begin{proposition}\label{prop:nonvanishing}
Let $f:(X,E)\to(Y,D)$ satisfy the hypotheses of
\Cref{thm:upper-bound}, and let
\[
  0\ne s\in H^0\bigl(X,m(K_X+E)\bigr),\qquad m>0.
\]
Then $\kappa(Y,L_Y)\ge0$, where $L_Y=K_Y+D$.
More precisely, apply the homogeneous highest-vector and marked-period
constructions to $s$, and denote their connected quotient by $S$.
If $\dim S>0$, then the relative Iitaka space $W$ of
\Cref{lem:relative-iitaka-construction,prop:rank-one-base} satisfies
\[
                 \kappa(Y,L_Y)\ge\dim W\ge\dim S>0.
\]
If $S$ is a point, there is an integer $a>0$ and a nonzero section
$\tau\in H^0(Y,aL_Y)$.  Moreover, for every $y\in V$ at which
$s|_{X_y}=0$, this section can be chosen with $\tau(y)=0$.
\end{proposition}

The construction is attached to the chosen section $s$; in particular,
its quotient $S$ may depend on $s$.  We first complete the point-quotient
case, using the original coefficient lattice.  The rest of the section
proves that a positive-dimensional quotient forces positive base
Kodaira dimension.

\subsection{The point-quotient case}

\begin{proof}[Proof of \Cref{prop:nonvanishing} when $S$ is a point]
For a point base there is a nonzero constant base section and the zero
test has no nonzero input satisfying its premise.  Assume $\dim Y>0$.
Apply \Cref{prop:logarithmic-vector} and
\Cref{prop:marked-period} to the chosen $s$, obtaining a homogeneous
vector with source $-aL_Y$.

\phantomsection\label{proof:point-nonvanishing}
Since $S$ is a point, the descended coefficient system is constant.  In
a constant frame its coordinates are rational sections of $aL_Y$.
The divisorial lattice bounds make each coordinate regular at every
prime of $Y$: this can be tested after finite ramified substitution,
since a negative order remains negative after multiplying it by the
ramification index.  On the smooth variety $Y$, regularity in
codimension one gives global regularity.  At least one coordinate
is nonzero because the homogeneous vector is nonzero; choose it as
$\tau$.

If $s|_{X_y}=0$ for $y\in V$, \Cref{lem:zero-test} says that all
these coordinates have positive order in the pulled-back original
$aL_Y$ frame along the divisor of the blowup of $y$ (along $y$ itself
if $Y$ is a curve).  Finite ramification multiplies orders by its
positive index.  Hence each original regular coefficient has positive
exceptional order, equal to its vanishing order at $y$.  Thus the same nonzero coordinate
vanishes at $y$, as asserted.
\end{proof}

\subsection{The numerical target}

Assume henceforth that $\dim S>0$.
After the birational preparations in
\Cref{lem:relative-iitaka-construction,prop:rank-one-base}, write
\[
  Y\xrightarrow{x}W\xrightarrow{h}S,\qquad
  u_*=\omega\,x^*u_W,
  \qquad L_W=K_W+T_W\sim_{\Q}K_W+B_W+M_W.
\]
Here $W,S$ are smooth and projective, both maps have geometrically
connected generic fibers, $0\le B_W\le1$ has SNC support, and $M_W$ is
nef.  The divisor $L_W$ is big over $S$.  The highest vector has source
$-aL_Y$, and $\omega$ is the generator in degree $a$ on the generic
$x$-fiber.  All references to $Y$ in this diagram use its prepared model;
the actual pluriform comparison in \Cref{prop:rank-one-base} returns
sections to the original pair.

\begin{proposition}[Removal of a period twist]\label{prop:twist-removal}
Let $h:W\to S$ be a surjective morphism of smooth projective
varieties, with $\dim S>0$. Let $L$ be a rational divisor on $W$
and $H$ a nef rational divisor on $S$. Suppose:
\begin{enumerate}
\item $L$ is big on the generic fiber of $h$, and
      $L-h^*K_S$ is pseudoeffective;
\item $K_S+bH$ is big for some rational $b\geq0$;
\item there are a dominant generically finite morphism
      $\pi:\widehat W\to W$ from a smooth projective variety,
      a nef rational divisor $P$ on $\widehat W$, and a rational
      $Q>0$ such that, with $\widetilde H=\pi^*h^*H$,
      \[
      \nu(P+\widetilde H)=\nu(P),\qquad Q\pi^*L\ge_{\mathrm{pe}} P.
      \]
\end{enumerate}
Then $L$ is big.
\end{proposition}

The first two hypotheses make $L+b h^*H$ big. The numerical-dimension
condition in the third says that the added period directions contribute
no positivity beyond $P$. Its fixed comparison with $L$ allows us to
remove the added divisor. We prove this numerical implication first;
the geometric comparisons that follow will establish its hypotheses
for $L_W$ and $\mathcal H$.

\begin{lemma}[Nef subtraction]\label{lem:nef-morse}
Let $Z$ be smooth projective of dimension $n>0$, let $A$ be an ample
rational divisor, and let $P,H$ be nef rational divisors satisfying
$\nu(P+H)=\nu(P)$.  For every fixed rational $b\ge0$, the divisor
$A+tP-bH$ is big for all sufficiently large rational $t$.
\end{lemma}

\begin{proof}
Put $r=\nu(P)$.  The polynomial
\[
                      (A+tP)^n
\]
has degree $r$ and strictly positive leading coefficient.  In the
polynomial $n b(A+tP)^{n-1}H$, all coefficients of $t^i$ for $i\ge r$
vanish.  Indeed
\[
 (P+H)^{i+1}A^{n-1-i}=0\quad(r\le i\le n-1),
\]
and each mixed nef intersection in this expansion is nonnegative;
in particular $P^iHA^{n-1-i}=0$.  When $r=n$, this index range is empty.
For $r=0$ this says that the entire
second polynomial is zero.  Thus, for large $t$,
\begin{equation}\label{eq:nonvanishing-morse}
                   (A+tP)^n>n(A+tP)^{n-1}bH.
\end{equation}

We recall directly why this strict nef Morse inequality implies
bigness here.  Write $U=A+tP$ and $V=bH$.  The first divisor is ample
and the second is nef.  Replace $V$ by $V+\delta A$, with rational
$\delta>0$ sufficiently small that the strict inequality persists.
After multiplying by a common positive integer, $U$ is integral ample
and the new $V$ is very ample.  Choose a general member $H_V\in|V|$.
Successive restriction exact sequences give
\[
 h^0\bigl(Z,l(U-V)\bigr)
   \ge h^0(Z,lU)
      -\sum_{j=0}^{l-1}
             h^0\bigl(H_V,(lU-jV)|_{H_V}\bigr).
\]
Multiplication by a nonzero section of $jV|_{H_V}$ bounds each
summand by $h^0(H_V,lU|_{H_V})$.  On a curve choose this section
nonvanishing at the finitely many points of $H_V$; in higher dimension
take $H_V$ smooth and irreducible.  Ample Hilbert asymptotics therefore
give
\[
 h^0\bigl(Z,l(U-V)\bigr)
 \ge\frac{U^n-nU^{n-1}V}{n!}\,l^n+O(l^{n-1}).
\]
The coefficient is positive.  Hence $U-V$ is big, and adding back
$\delta A$ proves bigness before perturbation as well.
\end{proof}

\begin{proof}[Proof of Proposition~\ref{prop:twist-removal}]
Choose $b\geq0$ with $K_S+bH$ big and put $D_b=L+b h^*H$.
By hypothesis \textup{(i)} and bigness of $K_S+bH$, there are an ample rational $A_S$ and
$\delta>0$ rational such that $D_b\ge_{\mathrm{pe}}\delta h^*A_S$.
The divisor $D_b$ is also big over $S$. Hence
$D_b+c h^*A_S$ is big for some rational $c>0$.

To justify the latter assertion, fix an ample $A_W$ on $W$.
Relative bigness gives a nonzero generic-fiber section of
$lD_b-A_W$ for some divisible $l$. Its direct image is nonzero;
after twisting by $kA_S$ for large $k$ it has a nonzero global
section. Thus $lD_b-A_W+k h^*A_S$ is effective. Taking $c=k/l$
proves the assertion. Now
\[
 D_b=\frac{\delta}{c+\delta}(D_b+c h^*A_S)
      +\frac{c}{c+\delta}(D_b-\delta h^*A_S)
\]
expresses $D_b$ as a positive multiple of a big class plus a
pseudoeffective class. Therefore $D_b$ is big.

Pull back and choose an ample rational $A'$ on $\widehat W$ with
$\pi^*L+b\widetilde H\ge_{\mathrm{pe}} A'$. For every $t>0$,
hypothesis \textup{(iii)} gives
\[
 (1+tQ)\pi^*L\ge_{\mathrm{pe}} A'+tP-b\widetilde H.
\]
The right side is big for large rational $t$ by
the numerical-dimension equality in \textup{(iii)} and
Lemma~\ref{lem:nef-morse}. Thus
$\pi^*L$, and hence $L$, is big. Bigness descends under a dominant
generically finite projective map, for example by the norm of a
section after subtracting a small ample class. No equality of
numerical and Iitaka dimensions for a general pseudoeffective
divisor is being assumed.

\end{proof}

\subsection{The fixed determinant lines and their cotangent maps}

Choose the connected unipotent level and equivariant resolutions of
\Cref{prop:tail-rank}, and write their composite as
$\pi:\widehat W\to W$.  This map is generically finite and Galois at the
function-field level.  On $\widehat W$ let $G_i$ be the extended subbundles
generated by contractions of $\theta^i(u_W)$, put $r_i=\rk G_i$, and retain
only indices with $r_i>0$.  Set
\begin{equation}\label{eq:nonvanishing-tail}
 P=-\sum_i(i+1)\det G_i,\qquad
 \widetilde{\mathcal H}=\pi^*h^*\mathcal H.
\end{equation}
Both lines are nef and
\begin{equation}\label{eq:nonvanishing-rank}
             \nu(P+\widetilde{\mathcal H})=\nu(P).
\end{equation}
The projective line $[u_W]$ gives precisely the compact-factor marks
required in that proposition, by the factorization
$u_*=\omega x^*u_W$.

Fix a rational frame of each $\det G_i^\vee$ over the original function
field $\C(W)$.  Its divisor in the extended line upstairs is denoted by
$\widehat D_i^G$.  Define rational divisors on $W$ by
\begin{equation}\label{eq:nonvanishing-determinants}
 D_i^G=\frac{\pi_*\widehat D_i^G}{\deg\pi},\qquad
 P_W=\sum_i(i+1)D_i^G.
\end{equation}
The superscript distinguishes these determinant divisors from the
boundary $D$ on $Y$.  The rational frames and these divisors are fixed
throughout the ample perturbations below.

We define the cotangent lattice before writing the transfer map. For
the rational SNC boundary $B_W=\sum_A b_AA$, its adapted cotangent
lattice $\Omega^1(W,B_W)$ has, at a general point of $A=(t=0)$,
transverse generator $t^{1-b_A}dt/t$ and ordinary tangential generators.
These are interpreted on a cover $t=v^e$ with $eb_A$ integral. Thus
$b_A=0$ retains the pullback of $dt$, and $b_A=1$ retains the pullback
of $dt/t$. An adapted cover is chosen with these divisible ramification
indices along the boundary; any auxiliary branch divisors have boundary
coefficient zero and retain the pullback of ordinary forms. Tensor maps
into this lattice are statements on such covers, in codimension one.
The proof will also verify their extension across the remaining locus.

\begin{lemma}[Orbifold transfer]\label{lem:orbifold-transfer}
The divisor $P_W$ is pseudoeffective and
\begin{equation}\label{eq:nonvanishing-pushdown}
                    \pi^*P_W\ge_{\mathrm{pe}}P.
\end{equation}
For each $i$, the determinant of the $i$th Higgs map gives a generically
nonzero rational map
\begin{equation}\label{eq:nonvanishing-orbifold-map}
 \OO_W(D_i^G-ar_iL_W)
       \longrightarrow \bigl(\Omega^1(W,B_W)\bigr)^{\otimes ir_i}.
\end{equation}
This notation means the following precise assertion: after a fixed
tensor power clearing the rational source divisor, the pullback map is
regular on every sufficiently divisible adapted cover, using the
orbifold cotangent lattice in codimension one.  The same tensor power
works after replacing $B_W$ by any effective SNC boundary
$B_W+\Lambda$ with coefficients at most one.
\end{lemma}

\begin{proof}
The extending lattices and their determinant divisors are invariant
under the Galois group.  For every prime $A_0$ of $W$, the orders of the
pullback of $D_i^G$ therefore agree with those of $\widehat D_i^G$ at all
primes over $A_0$.  Indeed, if $e$ is the common ramification index and
$v$ their common order, then the coefficient downstairs obtained by
pushforward and division by $\deg\pi$ is $v/e$.

Factor $\pi$ through the finite normalization $W^{\mathrm{fin}}\to W$.
The difference
\[
                    Z=\pi^*P_W-\sum_i(i+1)\widehat D_i^G
\]
is exceptional over $W^{\mathrm{fin}}$.  Its negative is relatively nef:
the pullback term has degree zero on contracted curves, and the remaining
term represents the nef line $P$.  The negativity lemma
\cite[Lemma~3.39]{KollarMori} gives $Z\ge0$.  This proves
\eqref{eq:nonvanishing-pushdown}.  Pushforward of the pseudoeffective
class $P$ gives $\deg(\pi)P_W$, proving the first assertion.

We prove the lattice assertion at an arbitrary prime $A_0\subset W$.
Use the notation of \Cref{prop:rank-one-base}:
\[
 t_{A_0}=\min_{I\to A_0}
       \frac{a^{-1}\operatorname{ord}_I(\omega)+d_I}{m_I},
 \quad m_I=\operatorname{ord}_I(x^*A_0),
 \quad d_I=\operatorname{coeff}_I(D).
\]
Choose a component $I$ achieving this minimum.  The discriminant
coefficient $b_{A_0}=\operatorname{coeff}_{A_0}(B_W)$ satisfies
\begin{equation}\label{eq:nonvanishing-same-minimizer}
 t_{A_0}=\frac{a^{-1}\operatorname{ord}_I(\omega)+d_I}{m_I},
 \qquad b_{A_0}\ge1-\frac{1-d_I}{m_I}.
\end{equation}
Both statements use this same $I$.

At general points of $A_0$ and $I$, choose a parameter $t$ downstairs
and a parameter $z$ upstairs.  After an etale local change absorbing a
unit, $t=z^{m_I}$; the other base coordinates pull back to independent
coordinates along $I$.  Work on a local cover $t=v^e$ sufficiently
divisible to make the Hodge monodromy unipotent and clear all relevant
denominators.  Write a coefficient $c$ of $\theta^i(u_W)$ in Schmid
frames, using an ordinary frame of $aK_W$ and a cotangent tensor having
$j$ factors $dt/t$, with the other factors tangential.  Define
$\operatorname{ord}(c)$ using the parameter $v$.

The identity $u_*=\omega x^*u_W$ holds for Higgs iterates as well:
the Higgs field is linear over functions, so there is no derivative of
$\omega$.  Make a further divisible substitution in $z$ dominating
the chosen $t$-cover.  The lattice bounds of
\Cref{prop:logarithmic-vector,lem:weight-extension} then apply to each
coefficient in the independent pulled-back coordinate tensors.
Each transverse factor pulls back as $m_I\,dz/z$ and costs
$1-d_I$ units of $I$-order.  Consequently
\[
 m_I\frac{\operatorname{ord}(c)}{e}
       +\operatorname{ord}_I(\omega)+ad_I-j(1-d_I)\ge0.
\]
The order of $\omega$ is measured in $aK_{Y/W}$, so the canonical
change of variables is already included.  Using
\eqref{eq:nonvanishing-same-minimizer}, we obtain
\begin{equation}\label{eq:nonvanishing-orbifold-order}
 \frac{\operatorname{ord}(c)}e
       \ge-at_{A_0}+\frac{j(1-d_I)}{m_I}
       \ge-at_{A_0}+j(1-b_{A_0}).
\end{equation}
This is exactly the required regularity after twisting the source by
$-aL_W$: the transverse orbifold cotangent frame is
$t^{1-b_{A_0}}dt/t$, interpreted after ramification, and the tangential
frames are ordinary differentials.  For example, when $b_{A_0}=0$
this frame is $dt$, and when $b_{A_0}=1$ it is $dt/t$.

Generically the contractions of this map surject onto $G_i$.  On the
unipotent resolution $G_i$ has its subbundle lattice inside the Schmid
grade, so the regularity just proved holds with this target lattice.
Dualize, take the top exterior power of $G_i^\vee$, and use the natural
inclusion of an exterior power in a tensor power.  The resulting source
is $\det G_i^\vee-ar_i\pi^*L_W$, which gives
\eqref{eq:nonvanishing-orbifold-map} by the order identification for
$D_i^G$ at primes over $W$.

Here is also a global interpretation of the cover assertion.  Fix an
integer $N_0>0$ clearing all the source divisors
$D_i^G-ar_iL_W$.  On a smooth adapted cover $\rho:W_{\rm ad}\to W$
for $B_W+\Lambda$, the $N_0$th tensor power is a rational map from
$\rho^*\OO_W(N_0(D_i^G-ar_iL_W))$ to the tensor power of degree
$N_0ir_i$ of its orbifold cotangent bundle.  At each prime it is regular
after a further common local ramification that is also unipotent for
the Hodge system.  Fractional cotangent frames pull back compatibly
under that ramification; an inequality of integral orders after such
a finite substitution implies the original inequality.  Thus the map
was already regular on $W_{\rm ad}$.  Regularity in codimension one
extends across codimension two because source and target are locally
free on its smooth model.

Increasing the boundary coefficient enlarges the cotangent lattice:
the original transverse frame is the new one multiplied by
$t^{\operatorname{coeff}_{A_0}(\Lambda)}$.  At auxiliary branch
divisors with zero boundary coefficient the lattice is the pullback
of ordinary cotangent forms.  Hence the same argument applies at every
prime of any adapted cover for the larger boundary.  Only the
denominators of this cover change with $\Lambda$; the integer $N_0$
is the fixed integer chosen above.
\end{proof}

\subsection{The uniform threshold and relative positivity}

We next turn these maps into a numerical comparison.  We use
\cite[Theorem~1.3]{CP} in its orbifold tensor form: for a smooth
projective SNC pair $(Z,\Delta)$ with rational coefficients in $[0,1]$
and $K_Z+\Delta$ pseudoeffective, every torsion-free quotient of an
orbifold cotangent tensor power on an adapted cover has nonnegative
degree against the pullback of every movable curve class on $Z$.
In particular, if a line injects into such a tensor power, its degree
is at most the degree of the determinant of that tensor power.  To
see the last assertion, saturate the image and apply the theorem to
the quotient; the effective divisor introduced by saturation only
strengthens the assertion.  For a rank-one tensor it follows directly
from the effective zero divisor of the map.

The threshold argument below follows the ample-perturbation method in
\cite[proof of Theorem~7.6]{CP}; here the fixed Higgs maps provide
constants independent of the varying boundary.

\begin{lemma}[A uniform pseudoeffective threshold]
\label{lem:pseff-threshold}
In the notation above, $L_W$ is pseudoeffective.  There is a constant
$Q>0$, depending only on the fixed Higgs maps, such that
\begin{equation}\label{eq:nonvanishing-domination}
                         Q\pi^*L_W\ge_{\mathrm{pe}}P.
\end{equation}
\end{lemma}

\begin{proof}
Fix an ample rational divisor $A$ on $W$ and write $n=\dim W>0$.
For a positive rational $\epsilon$ for which $L_W+\epsilon A$ is
pseudoeffective, the divisor $M_W+\epsilon A$ is ample.  Choose a
sufficiently divisible $k$ and a general smooth member of
$|k(M_W+\epsilon A)|$, transverse to every stratum of $B_W$, and divide
it by $k$.  This gives
\[
  \Lambda_\epsilon\sim_{\Q}M_W+\epsilon A,
  \quad 0\le B_W+\Lambda_\epsilon\le1,
  \quad K_W+B_W+\Lambda_\epsilon\sim_{\Q}L_W+\epsilon A.
\]
The support is SNC.  Thus the pair is log canonical and has
pseudoeffective adjoint, as required in \cite[Theorem~1.3]{CP}.

Choose a smooth Galois adapted cover for this pair, as in
\cite[Definition~5.1 and Lemma~5.2]{CP}.  Such covers can be
obtained by the covering construction: add general transverse
auxiliary divisors to make the root divisors divisible in the Picard
group, and normalize the resulting finite root covers.  Use sufficiently
many general auxiliary choices that, at each point and for each
prescribed boundary component, one auxiliary equation is a unit.
Locally the normalization then extracts roots of independent SNC
coordinates and is smooth.  The auxiliary ramifications are assigned
boundary coefficient zero.  On this cover the determinant of the
orbifold cotangent bundle is the pullback of
$K_W+B_W+\Lambda_\epsilon$.

Put $d_i=ir_i$.  For $d_i>0$, the tensor in
\Cref{lem:orbifold-transfer} has rank $n^{N_0d_i}$ and first Chern
class
\[
 N_0d_i\,n^{N_0d_i-1}
             \rho^*(L_W+\epsilon A).
\]
Apply the preceding quotient statement to its nonzero source line.
Divide the resulting inequality by $N_0\deg\rho$.  For every movable
curve class $\alpha$ on $W$ this gives
\[
 \bigl(c_i(L_W+\epsilon A)+ar_iL_W-D_i^G\bigr)\cdot\alpha\ge0,
 \qquad c_i=d_i n^{N_0d_i-1}.
\]
For $d_i=0$, the map is into the trivial line; take $c_i=0$ and use
its effective zero divisor.  All $c_i$ are independent of $\epsilon$
and of the chosen representative $\Lambda_\epsilon$.

The duality between the pseudoeffective divisor cone and the movable
curve cone on the smooth projective variety $W$
\cite[Theorem~2.2]{BDPP} now gives
\begin{equation}\label{eq:nonvanishing-uniform-inequality}
 C(L_W+\epsilon A)+RL_W\ge_{\mathrm{pe}}P_W,
 \qquad
 C=\sum_i(i+1)c_i\ge0,
 \quad R=a\sum_i(i+1)r_i>0.
\end{equation}
The strict positivity of $R$ follows from $G_0\ne0$.

For completeness, consider
\[
 \tau=\inf\{t\ge0:L_W+tA\text{ is pseudoeffective}\}.
\]
This is finite.  Since $P_W$ is pseudoeffective,
\eqref{eq:nonvanishing-uniform-inequality} implies, for every rational
$\epsilon>\tau$,
\[
        L_W+\frac{C}{C+R}\epsilon A
                         \text{ is pseudoeffective}.
\]
If $\tau>0$, closedness of the cone and $\epsilon\downarrow\tau$
would give $\tau\le C\tau/(C+R)<\tau$, a contradiction.  Thus
$\tau=0$ and $L_W$ is pseudoeffective.  We may now let
$\epsilon\downarrow0$ in
\eqref{eq:nonvanishing-uniform-inequality}; pullback and
\eqref{eq:nonvanishing-pushdown} yield
\eqref{eq:nonvanishing-domination} with $Q=C+R$.
\end{proof}

\Needspace{8\baselineskip}
\begin{lemma}[Relative pseudoeffectivity]\label{lem:nonvanishing-relative}
For the prepared morphism $h:W\to S$ one has
\begin{equation}\label{eq:nonvanishing-relative}
                         L_W-h^*K_S\ge_{\mathrm{pe}}0.
\end{equation}
\end{lemma}

\begin{proof}
If $h$ has relative dimension zero, it is birational by connectedness
of its generic fiber, and the conclusion follows from
$K_W-h^*K_S\ge0$, $B_W\ge0$, and nefness of $M_W$.
Assume henceforth that its relative dimension is positive.
Use $A$ and $\Lambda_\epsilon$ from the preceding proof.  We have
established pseudoeffectivity of
$K_W+B_W+\Lambda_\epsilon\sim_{\Q}L_W+\epsilon A$ for every rational
$\epsilon>0$.  Apply \cite[Theorem~3.4 and Remark~3.3]{CP} to the
fibration and this smooth SNC log canonical pair.  We spell out the
model issue because the fixed $W$ need not itself be the good model
used in that theorem.

Take a commutative birational diagram
\[
 \begin{array}{ccc}
  W'&\xrightarrow{h'}&S'\\
  {\scriptstyle p}\downarrow&&\downarrow{\scriptstyle q}\\
  W&\xrightarrow{h}&S,
 \end{array}
\]
with smooth projective varieties, resolved morphism, SNC discriminant
and reduced inverse discriminant, and such that every divisor of $W'$
contracted by $h'$ to codimension at least two is exceptional over
$W$.  These are the good-model requirements in
\cite[Remark~3.3]{CP}.  They may be arranged by extracting on $S$ the
restricted valuations of the finitely many divisors of $W$ with
exceptional image, as in \Cref{lem:valuation-preparation}, and then
resolving the map and the boundary.  Once those original valuations
have divisorial center, further base modifications preserve that
property for their strict transforms; all remaining offending source
divisors are exceptional over $W$.

One can require the discriminant here to be the actual critical-value
locus, rather than merely an SNC divisor containing it.  To check this
point in the preparation, first take an SNC divisor on the resolved
base containing the critical values and all the loci over which the
modifications occur; resolve its reduced inverse image, preserving
the smooth complement.  If a component of this base divisor is not
generically a critical-value divisor and the relative dimension is
positive, blow up a smooth component of its inverse image intersected
with a general very ample hypersurface.  Choose a component of this
intersection dominating the base divisor, which exists since its
generic fiber has positive dimension.  The center has normal crossings
with the other boundary components.  Locally it has equations
$t=z=0$, where $t$ is the base parameter and $z$ a fiber coordinate;
after the blowup a chart has $t=uv$.  The resulting critical locus
therefore dominates that base divisor.  Repeating for the finitely
many remaining components makes the actual critical-value locus
equal to the chosen SNC divisor.  All these new source divisors are
exceptional over $W$, and the reduced inverse image remains SNC.

Set $\Delta_\epsilon=B_W+\Lambda_\epsilon$ and take on $W'$ the
strict transform of $\Delta_\epsilon$ plus the reduced
$p$-exceptional divisor, resolving its support together with the
discriminant.  Denote this boundary by $\Delta'_\epsilon$.  Log
canonicity gives
\[
 K_{W'}+\Delta'_\epsilon
     =p^*(K_W+\Delta_\epsilon)+E_\epsilon,
               \qquad E_\epsilon\ge0
               \text{ exceptional over }W.
\]
Its adjoint is therefore pseudoeffective.  Theorem~3.4 of \cite{CP}
is applied with the full boundary $\Delta'_\epsilon$ and gives
pseudoeffectivity of
\[
 K_{W'/S'}+(\Delta'_\epsilon)_{\rm hor}-D(h'),
\]
where $D(h')$ counts ramification along primes dominating divisors of
$S'$.  Its generic-fiber hypothesis, as used in the proof of that
theorem, holds because the restriction of the pseudoeffective total
adjoint to a very general fiber is pseudoeffective.  The vertical
boundary restricts to zero there, so the fiber sees exactly
$(\Delta'_\epsilon)_{\rm hor}$.  This does not require
$K_{W'}+(\Delta'_\epsilon)_{\rm hor}$ to be pseudoeffective globally;
see also \cite[Remark~3.6]{CP}.
Adding the effective vertical boundary and $D(h')$ proves that
$T'_\epsilon=K_{W'/S'}+\Delta'_\epsilon$ is pseudoeffective.

Write $K_{S'}=q^*K_S+F$ with $F\ge0$ exceptional; this is the
canonical Jacobian divisor of a birational morphism between smooth
varieties.  Pushforward of the preceding pseudoeffective divisor gives
\[
 p_*T'_\epsilon
   =K_W+\Delta_\epsilon-h^*K_S-p_*h'^*F.
\]
Consequently $L_W+\epsilon A-h^*K_S$ is the sum of a
pseudoeffective class and the effective class $p_*h'^*F$.
Letting $\epsilon\downarrow0$ proves
\eqref{eq:nonvanishing-relative} on the original fixed $W$.
\end{proof}

\subsection{Return to the original base}

The geometric comparisons now supply all three hypotheses of the
numerical removal proposition.  Its conclusion gives a big divisor
on $W$; the exact section lattice then returns actual pluriforms to
the original base.

\begin{proposition}\label{prop:positive-period}
For the data attached to a nonzero total-space pluriform, if
$\dim S>0$ then $L_W$ is big and
\[
                      \kappa(Y,L_Y)\ge\dim W>0.
\]
\end{proposition}

\begin{proof}
Apply \Cref{prop:twist-removal} with $L=L_W$ and
$H=\mathcal H$. Relative bigness is
\Cref{prop:rank-one-base}; relative pseudoeffectivity is
\Cref{lem:nonvanishing-relative}; adjoint bigness is
\Cref{prop:adjoint-bigness}. The nef line $P$ and numerical-rank equality
are supplied by \Cref{prop:tail-rank}; its fixed comparison with $L_W$
is \Cref{lem:pseff-threshold}. Thus $L_W$ is big.

The divisible-degree comparison of \Cref{prop:rank-one-base} sends its
complete section spaces to actual sections on the original fixed base
and preserves their ratios. The images of sufficiently large divisible
systems therefore have dimension $\dim W$. Hence
$\kappa(Y,L_Y)\geq\dim W\geq\dim S>0$.
\end{proof}

\begin{proof}[Proof of \Cref{prop:nonvanishing}]
The point-quotient case, including its original-point zero assertion,
was proved at the start of the section.  For $\dim S>0$,
\Cref{lem:relative-iitaka-construction} supplies the nonnegative
generic-fiber Iitaka dimension and the diagram used above.
\Cref{prop:positive-period} then gives
$\kappa(Y,L_Y)\ge\dim W\ge\dim S>0$.  These two alternatives prove
all the assertions.
\end{proof}

\begin{proof}[Proof of \Cref{prop:section-construction}]
Apply \Cref{prop:nonvanishing} to the specified nonzero section $s$.
It gives $\kappa(Y,L_Y)\geq0$. If that dimension is zero, the
positive-dimensional quotient case is excluded. The point-quotient
case supplies a nonzero original base section vanishing at every
chosen $y\in V$ for which $s|_{X_y}=0$. This is the second assertion.
Together with the reduction in \Cref{sec:completion}, it completes the
independent upper inequality.
\end{proof}

\subsection{Equality from logarithmic subadditivity}
\label{sec:matched-equality}

We finally invoke the only companion input used in this paper. Its
statement is recorded here to make the equality's precise hypothesis
match explicit.

\begin{theorem}[Logarithmic subadditivity, {\cite[Corollary~\FoundationLogNumber]{Foundation}}]
\label{thm:log-iitaka}
Let $g\colon Z\to T$ be a surjective morphism with connected fibers between
smooth connected projective complex varieties. Let $B_Z,B_T$ be reduced
SNC divisors with
$\Supp(g^*B_T)\subseteq\Supp(B_Z)$, and let
$G$ be a very general smooth fiber with $B_G=B_Z|_G$. Then
\begin{equation}\label{eq:lower-bound}
 \kappa(Z,K_Z+B_Z)
 \geq \kappa(T,K_T+B_T)+\kappa(G,K_G+B_G).
\end{equation}
No smoothness outside $B_T$, abundance, or good-minimal-model hypothesis
is imposed.
\end{theorem}

To prove \Cref{cor:additivity}, apply logarithmic subadditivity
\cite[Corollary~\FoundationLogNumber]{Foundation}, stated in \Cref{thm:log-iitaka}, with
$Z=X$, $T=Y$, $g=f$, $B_Z=E$, and $B_T=D$. The source and base are smooth
connected projective complex varieties, the morphism is surjective with
connected fibers, the boundaries are reduced SNC, and
\eqref{eq:support-inclusion} is exactly its boundary hypothesis.
The very general smooth fiber and its restricted boundary are precisely
$F,E_F$. This theorem supplies the lower inequality; \Cref{thm:upper-bound}
supplies the upper inequality. Their infinity conventions agree, giving
\eqref{eq:additivity} in every case.
